% Section 9.2, from lectures/L09/appendix/b_time.md.

\section{Time, Delays and Timers}
\label{c9:app:b}

Four ways to say ``later'', how much later each of them actually means, and why the difference is
larger than anyone expects.

Every number here was measured on this course's target with the lab's \code{qa_timing} module, at
\code{HZ=250}. Yours will match, because they are properties of the configuration rather than of the
machine.

\subsection{Jiffies, and why not to quote them}
\label{c9:app:b1}

The kernel counts ticks in a global \code{jiffies}. The tick rate is \code{HZ}, a compile-time
constant:

\begin{console}
# zcat /proc/config.gz | grep '^CONFIG_HZ='
CONFIG_HZ=250
\end{console}

\noindent so on this kernel \textbf{one jiffy is 4 milliseconds}. Common values are 100, 250, 300
and 1000, and which one a distribution picks is a trade between timer resolution and tick overhead.

\textbf{A jiffy is not a unit of time.} It is a unit of ticks, and the same code on a \code{HZ=1000}
kernel waits a quarter as long. Any duration you care about should be written in real units and
converted:

\begin{ccode}
unsigned long timeout = jiffies + msecs_to_jiffies(500);
if (time_after(jiffies, timeout)) { /* ... */ }
\end{ccode}

\noindent \code{time_after} and \code{time_before} rather than \code{<} and \code{>}, because
\code{jiffies} wraps. On a 32-bit counter at \code{HZ=250} that is once every 199 days, and the
comparison macros handle it; a naive comparison produces a driver that misbehaves twice a year.

\subsection{\code{ktime}, for measuring}
\label{c9:app:b2}

Jiffies are for \emph{scheduling}. For \emph{measuring}, use \code{ktime}, which is nanoseconds:

\begin{ccode}
ktime_t start = ktime_get();
/* ... */
s64 us = ktime_us_delta(ktime_get(), start);
\end{ccode}

\begin{booktable}{@{}>{\hsize=0.74\hsize}L>{\hsize=1.11\hsize}L>{\hsize=1.15\hsize}L@{}}
  \thead{Clock} & \thead{Reads} & \thead{Use for} \\
  \midrule
  \code{ktime_get} & Monotonic since boot & \textbf{Intervals.} Always this one \\
  \code{ktime_get_real} & Wall-clock time & Timestamps a human will read \\
  \code{ktime_get_boottime} & Monotonic, includes suspend & Intervals across a sleep \\
\end{booktable}

\noindent \textbf{Measure intervals with the monotonic clock.} Wall-clock time can jump forwards or
backwards when NTP corrects it or somebody sets the date, and an interval computed across that is
nonsense, occasionally negative. This is the same rule userspace has with \code{CLOCK_MONOTONIC},
and it is the one Chapter~\ref{c12:ch} depends on entirely.

\subsection{Delaying against sleeping}
\label{c9:app:b3}

Two families, and confusing them is a correctness bug rather than a performance one.

\begin{booktable}{@{}>{\hsize=0.87\hsize}LlL>{\hsize=1.11\hsize}L@{}}
  \thead{Function} & \thead{Mechanism} & \thead{Context} & \thead{Costs} \\
  \midrule
  \code{udelay(n)} & Busy-wait & \textbf{Any}, including atomic & A CPU, entirely \\
  \code{mdelay(n)} & Busy-wait & Any & A CPU, for milliseconds \\
  \code{msleep(n)} & Timer wheel & Process context only & Nothing \\
  \code{usleep_range(a, b)} & hrtimer & Process context only & Nothing \\
\end{booktable}

\noindent \textbf{The \code{delay} family burns a CPU} and is legal in atomic context precisely
because it never sleeps. \textbf{The \code{sleep} family gives the CPU up} and cannot be used where
sleeping is forbidden.

The rule: \textbf{use \code{udelay} only for short hardware delays in code that cannot sleep}, and
only for a few microseconds. \code{udelay(5000)} is legal, compiles, works, and spends five
milliseconds of a CPU doing nothing, which is a real-time problem that Chapter~\ref{c12:ch} will
come back to. \code{mdelay} exists mostly so that its name can warn you.

\subsection{What the sleeps actually cost}
\label{c9:app:b4}

This is the measurement that surprises people. From the lab's \code{qa_timing}, 200 repetitions of
each, at \code{HZ=250} where one jiffy is 4000 us:

\begin{narrowtable}{@{}lrrr@{}}
  \thead{Call} & \thead{Asked for} & \thead{Measured, each} & \thead{Ratio} \\
  \midrule
  \code{msleep(1)} & 1,000 us & \textbf{8,000 us} & 8.0$\times$ \\
  \code{msleep(4)} & 4,000 us & \textbf{8,000 us} & 2.0$\times$ \\
  \code{msleep(10)} & 10,000 us & 15,998 us & 1.6$\times$ \\
  \code{usleep_range(1000)} & 1,000 us & 1,443 us & 1.4$\times$ \\
  \code{usleep_range(100)} & 100 us & 307 us & 3.1$\times$ \\
  \code{udelay(100)} & 100 us & 112 us & 1.1$\times$ \\
\end{narrowtable}

\noindent \textbf{\code{msleep(1)} and \code{msleep(4)} take exactly the same time.} Asking for a
quarter as long got nothing. That single row is the most useful thing in this section:
\code{msleep} is quantised to jiffies, and below one jiffy the argument stops carrying information.

The arithmetic is exact and you can reproduce it. \code{msecs_to_jiffies} rounds \textbf{up}:

\begin{narrowtable}{@{}llll@{}}
  \thead{Call} & \thead{\code{msecs_to_jiffies}} & \thead{Plus the guarantee} & \thead{Total} \\
  \midrule
  \code{msleep(1)} & 1 jiffy & $+1$ & 2 jiffies = 8 ms \\
  \code{msleep(4)} & 1 jiffy & $+1$ & 2 jiffies = 8 ms \\
  \code{msleep(10)} & 3 jiffies & $+1$ & 4 jiffies = 16 ms \\
\end{narrowtable}

\noindent Where the extra jiffy comes from is worth knowing, because it is not in \code{msleep}'s
source. \code{schedule_timeout}'s own documentation says:

\begin{quote}
at least \code{@timeout} jiffies are guaranteed to pass before the routine returns
\end{quote}

\noindent\textbf{``At least'' is the whole answer.} You call it partway through a jiffy, and the kernel does
not know how far through. To guarantee that N jiffies have \emph{elapsed}, it must wait for N+1
ticks. The overshoot is the price of a guarantee that never returns early.

\textbf{\code{usleep_range} is the fix}, and the table shows why: it uses hrtimers rather than the
jiffy wheel, so a 1 ms request costs 1.4 ms instead of 8 ms. It takes a range because giving the
kernel slack lets it coalesce your timer with others already due, which saves wakeups on a system
trying to stay idle. Use it for anything between about 10 us and 20 ms; \code{msleep} above that,
where the jiffy quantum stops mattering.

\lecturefigure[0.78]{lectures/L09/appendix/images/sleep_cost.png}
  {Horizontal bars comparing what each call asked for against what it measured, on a logarithmic axis,
   with dashed lines at two and four jiffies. \code{msleep(1)} and \code{msleep(4)} both land on 8,000
   microseconds, and \code{msleep(10)} on 15,998.}
  {fig:sleep_cost}

\subsection{Timers}
\label{c9:app:b5}

Two kinds, and the difference is the same jiffy-against-nanosecond split.

\subsubsection{\code{timer_list}, on the jiffy wheel}
\label{c9:app:b:timer_list-on-the-jiffy-wheel}

\begin{ccode}
static struct timer_list my_timer;

timer_setup(&my_timer, my_callback, 0);
mod_timer(&my_timer, jiffies + msecs_to_jiffies(500));
\end{ccode}

\noindent Resolution is a jiffy, and the wheel deliberately batches timers into buckets, so a long
timeout may be rounded by a noticeable fraction. That is a feature: it lets an idle system stay
idle.

\textbf{The callback runs in softirq context}, so it may not sleep, may not take a mutex, and may
not call \code{copy_to_user}. This catches people, because the code around a timer is usually
process context and the callback looks like it belongs to it.

\code{timer_delete_sync} before freeing anything the callback touches, and note the \emph{sync}: it
waits for a callback already running on another CPU. \code{timer_delete} alone can return while the
callback is still executing on the memory you are about to free. Most code you will read calls them
\code{del_timer_sync} and \code{del_timer}, the old names, which 6.12 keeps as wrappers marked not
for new code.

\subsubsection{\code{hrtimer}, on real time}
\label{c9:app:b:hrtimer-on-real-time}

\begin{ccode}
hrtimer_init(&my_hrtimer, CLOCK_MONOTONIC, HRTIMER_MODE_REL);
my_hrtimer.function = my_callback;
hrtimer_start(&my_hrtimer, ms_to_ktime(500), HRTIMER_MODE_REL);
\end{ccode}

\noindent Note that the callback is assigned to the structure rather than passed to the initialiser.
That is the 6.12 interface; a later kernel introduces \code{hrtimer_setup}, which takes the callback
as an argument and is what newer material will show you. This is \secref{c3:app:a2}'s unstable
internal API again, in a function you would meet in your first week.

Nanosecond resolution, backed by a real hardware timer rather than the tick. The callback returns
\code{HRTIMER_NORESTART} or \code{HRTIMER_RESTART}, which makes periodic timers straightforward.

hrtimers cost more per timer, so the kernel uses the wheel for the many timeouts that never expire
(network retransmits, mostly) and hrtimers for the few that must be accurate.

\subsubsection{Choosing}
\label{c9:app:b:choosing}

\begin{booktable}{@{}>{\hsize=0.72\hsize}L>{\hsize=1.28\hsize}L@{}}
  \thead{Use} & \thead{Because} \\
  \midrule
  \code{timer_list} for a timeout & It will usually be cancelled, and precision is irrelevant \\
  \code{hrtimer} for a deadline & You need the resolution, and it will fire \\
  Neither, if an interrupt will do & The device already knows when it has data \\
\end{booktable}

\noindent That last row is worth stating. This course's driver could poll the device on a timer; it
does not, because Chapter~\ref{c8:ch} gave it an interrupt. \textbf{A timer is what you use when
nothing tells you.} The exercises ask you to build the timer version anyway and compare, because the
comparison is the argument.

\subsection{What none of this gives you}
\label{c9:app:b6}

\textbf{None of it is a guarantee about the maximum.} Everything above is ``at least'', and the
numbers in \secref{c9:app:b4} are averages over 200 repetitions on an idle machine. What happens to
the worst case under load is not in this section, is not what \code{msleep} documents, and is the
subject of Chapter~\ref{c12:ch}.

\textbf{And the measurements themselves contain the emulator.} The 1,443 us for a 1 ms
\code{usleep_range} is this target's number, and part of it belongs to QEMU rather than to Linux.
That is the same caution the Cross-check of Chapter~\ref{c8:ch}, Exercise~\ref{c8:ex:9}, arrives at,
and Chapter~\ref{c12:ch} states it in full.
